\documentclass[11pt,a4paper]{article}
\usepackage[T1]{fontenc}
\usepackage[utf8]{inputenc}
\usepackage{lmodern,amsmath,amssymb,booktabs,array,xcolor,microtype}
\usepackage[a4paper,left=2.3cm,right=2.3cm,top=2.5cm,bottom=2.5cm]{geometry}
\usepackage[colorlinks=true]{hyperref}
\definecolor{ojacblue}{rgb}{0,0,0.8}
\hypersetup{linkcolor=ojacblue,citecolor=ojacblue,urlcolor=ojacblue}
\renewcommand{\baselinestretch}{1.08}
\setlength{\parindent}{0pt}
\setlength{\parskip}{4pt}
\setlength{\emergencystretch}{3em}
\newcommand{\R}{https://github.com/ralfanop/Certifications_Eulerian_polar-curve_computations}
\newcommand{\Br}{gegenbauer-fock-formalization-20261025}
\newcommand{\eq}[1]{(#1)}
\newcommand{\rE}{r_E}
\begin{document}

\begin{center}
{\Large\bfseries Manuscript 13: Lean~4 and Arb certification record}\\[6pt]
{\large Eulerian Fourier profiles and Gegenbauer--Fock operators: dimensional and spectral bridges}\\[6pt]
Roberto Alfano, 7~November~2026
\end{center}

\section*{1. Scope.}
This record states, for manuscript~13, what is kernel-checked in Lean~4 and which numerical values are certified in exact rational and Arb ball arithmetic. It replaces Appendix~B (Lean~4 certification) of the article and the description of the certificates formerly given in its Code Availability statement. Equation, section and result numbers are those of the version of 2~November~2026 with the notation correction of 5~November~2026, in which the Hilbert spaces are written $\mathcal H_\lambda$.

\section*{2. The Lean~4 library.}
The library \texttt{GegenbauerFock}, in the directory \texttt{lean/GegenbauerFock} of the repository \url{\R} on the branch \href{\R/tree/\Br}{\texttt{\Br}}, is compiled with the Lean~4 toolchain \texttt{v4.35.0-rc3} and Mathlib commit \texttt{0272bbe1c093}. On 2~November~2026 (commit \href{\R/tree/1132049025418c06f8b6f16a8e0112ab3088c4ba}{\texttt{11320490}}) its 792 modules contained 2547 theorems and lemmas. Continuous integration recompiles the modules below from source and prints the axioms of their main statements; the audit found only \texttt{propext}, \texttt{Classical.choice} and \texttt{Quot.sound}, and no \texttt{sorry}, \texttt{admit}, added axiom or \texttt{native\_decide}.

Every numbered statement of the article is formalised:
\begin{itemize}\setlength{\itemsep}{1pt}\raggedright
\item Definition~2.1 and Proposition~3.1 in \texttt{Manuscript13Measure};
\item Proposition~3.2 by the identification of the image of the beta law under $u=2x-1$ with $\nu_\lambda$;
\item Propositions~3.3, 3.4 and~3.6 in \texttt{Section3Ultraspherical};
\item Section~4 and Table~1 in \texttt{Section4Resolvent};
\item Proposition~5.1 and Corollary~5.2 in the Sturm--Liouville modules;
\item Definition~6.1 and Theorems~6.2--6.4 in the ultraspherical Fock modules and \texttt{Section6Corollaries};
\item Propositions~6.5 and~6.6 in \texttt{Section6Corollaries} and \texttt{EulerianStateShi};
\item Section~7 in \texttt{Section7Intertwining} and \texttt{Section7Closed};
\item Theorem~8.1, Propositions~8.2 and~9.1 and Lemma~9.2 in \texttt{Section8Section9};
\item Proposition~10.1 and the identities of Section~10 in \texttt{Section10Comparison}.
\end{itemize}

\section*{3. Classical inputs.}
Classical inputs enter as explicit hypotheses and are not results of the article: the reality of the Neumann function $Y_\nu$ and the closed form of $Y_{1/2}$, used only for the complete kernels of Section~4, the zonal marginal law behind \eq{38}, and the Funk--Hecke formula behind Proposition~10.1. The closed forms of $J_{\pm1/2}$ are proved.

\section*{4. Formalisation conventions.}
Mathlib has no singular values, Schatten classes or continuous-spectrum predicate, so \eq{72} is formalised as the existence of an operator of rank at most $N$ within $\exp(c/2)(c^N/N!)^{1/2}$ of $C_{\lambda,\alpha}$ in operator norm, Schatten membership as the summability of these bounds raised to every power $p>0$, and the last clause of Proposition~8.2 as $\ker S_{\lambda,c}=\{0\}$ with dense range and without surjectivity. The complex operator $D_\lambda$ of Theorem~7.3 is the complexification of the real selfadjoint realisation of the library, and the values at $\rho=0$ in Section~4 are stated as limits.

\section*{5. Numerical certificates.}
The script \href{\R/blob/main/numerics/arb/v007_arb.py}{\texttt{numerics/arb/v007\_arb.py}}, with its log \texttt{v007\_arb.log}, uses exact rational arithmetic and Arb balls at 256 bits (\texttt{python-flint}); every comparison is exact or a certain Arb comparison. It certifies: for Proposition~6.5, the fifteen-term partial sum of $I=\sum_n1/[4^nn!(2n+1)]$, its remainder bound with $e^{1/4}<4/3$ and the printed decimal enclosure, and, independently, that $I=\sqrt\pi\operatorname{erfi}(1/2)$ lies inside that enclosure; for Table~1, $A_d(k)k^{d-2}=4\Gamma(d/2)(4\pi)^{(d-2)/2}$ for $d=3,\ldots,12$ and the endpoint values $2k$ ($d=1$) and $4$ ($d=2$), as exact rational multiples of powers of $\pi$; for Corollary~4.2, the factor $8\pi$; for Proposition~6.6 at $\alpha=1$, the agreement of $\sum_k1/[16^k(2k)!(2k+1)^2]$ with $4\operatorname{Shi}(1/4)$.

The directory \href{\R/tree/main/figures/13_gegenbauer_fock}{\texttt{figures/13\_gegenbauer\_fock}} holds the sources, data and scripts of the figures, and the certificates of those whose captions state numerical facts. For Figure~4(b), \texttt{cert\_fig3\_gram.py} computes the normalised Gram matrix $\langle F_m,F_n\rangle/(\|F_m\|\|F_n\|)$, $0\leq m,n\leq11$, of $F_n=C_{1/2,\alpha}P_n$ with $\alpha=1/16$ ($c=4$) from exact rational partial sums up to $K=90$, with tail at most $55c^K/K!<6\cdot10^{-83}$, and certifies that the entries with $m+n$ odd vanish, that those with $m+n$ even are positive and decrease away from the diagonal, and that the largest off-diagonal entry, at $(m,n)=(0,2)$, is $0.753174039953$ to twelve digits. For Figure~5, \texttt{make\_fig1\_fig4\_data.py} encloses the tail $\|C-C_N\|_{\mathrm{HS}}^2=\sum_{n\geq N}c^n/[(2n+1)n!]$ for $\lambda=1/2$, $\alpha\in\{1/4,1,4\}$ and $0\leq N\leq20$ by exact rational partial sums and a rigorous remainder, and checks the inequality $\|C-C_N\|_{\mathrm{HS}}^2\leq e^cc^N/N!$ of Theorem~8.1 exactly at all 63 points.

\end{document}
